\section{K-016: removing the parity restriction in the single-pair criterion}

\begin{theorem}[K-016]
Let $n\ge4$, without a parity restriction.  Let $T$ be a concise
tensor with $L^0=S$, and suppose $g\in L$ has order $M=2^m$,
primitive determinant character of order $M$, and
$C_L(g)=S\langle g\rangle$.  Suppose that in each factor the
number of antipodal eigenvalue-index pairs is at most one, in the
precise sense used in K-013.  If a factor has no antipodal pair, or
if one of the three one-paired-subspace projections of K-013 is
nonzero, then
\[
                         L=S\langle g\rangle.
\]
In particular the family of K-015 has full stabilizer
$S\times\mu_{2^{n-1}}$ in every dimension $n\ge5$.
\end{theorem}

\begin{proof}
The only parity-dependent steps in K-012 were the descent of the
centralizer and the choice of an order-$M$ representative in the
determinant-normalized central extension.  We replace both, and
then give the character argument with the revised order.

Each factor has an unpaired eigenvalue, since $n\ge4$ and there
is at most one antipodal pair.  A multiset invariant under
multiplication by a nontrivial two-power root is invariant under
multiplication by $-1$: the latter lies in the cyclic subgroup
generated by that root.  Therefore none of the factor spectra has
such a nontrivial scalar symmetry.  If $hgh^{-1}=sg$ with $s\in S$,
taking $M$th powers gives $s_i^M=1$ in each factor.  Conjugacy
preserves its spectrum, so the preceding observation gives
$s_i=1$.  Hence
\[
                 C_{L/S}(\bar g)=\langle\bar g\rangle.
                                                               \tag{K-016a}
\]
The column-orthogonality argument of K-012 consequently still gives
$H^{\rm ab}\simeq\mu_M$ via $\chi$, where $H=L/S$.

Again set
$\widetilde H=\{(A,B,C)\in L:\det A=\det B=1\}$,
with central kernel $Z\simeq\mu_n^2$.  Write $n=2^e n_0$ with
$n_0$ odd.  Raising to the $n$th power maps
$\mu_{M2^e}$ onto $\mu_M$.  We can therefore multiply the three
factors of $g$ by two-power scalar roots whose product is one and
obtain $g_0\in\widetilde H$.  Its order $Q$ is a power of two
dividing $M2^e$, and its image $\bar g$ still has order $M$.
The scalar normalization preserves all eigenvalue ratios, all
antipodal pairs, and all eigenspaces.

For every odd $u$,
\[
 C_{\widetilde H}(g_0^u)=Z\langle g_0\rangle,
               \qquad |C_{\widetilde H}(g_0^u)|=|Z|M.
                                                               \tag{K-016b}
\]
Indeed its image in $H$ is contained in the cyclic centralizer
from (K-016a), and every element of that preimage is $zg_0^r$ and
commutes.  If $z\in Z$ is nonidentity, $g_0^u$ and $g_0^uz$ are
not conjugate.  To see this, their factorwise spectra would have
to be invariant under multiplication by the corresponding $z_i$.
All eigenvalue ratios in each factor of $g_0^u$ lie in $\mu_M$,
so such a scalar $z_i$ must lie in $\mu_M$.  If it is nontrivial,
the factor spectrum would be entirely antipodally paired,
contrary to the first paragraph.  Thus conjugacy would force
$z_i=1$ in all factors and $z=1$.

Fourier inversion of character column orthogonality on $Z$ now
gives, for every central character $\omega$ and every odd $u$,
\[
 \sum_{\substack{\rho\in\operatorname{Irr}(\widetilde H)\\
                    \rho|_Z=(\dim\rho)\omega}}
                             |\rho(g_0^u)|^2=M.             \tag{K-016c}
\]
The right side is $M$, not the potentially larger order $Q$.
An irreducible constituent with no nontrivial $\chi$ self-twist
has $M$ distinct twists in this central-character block.  Thus
$|\rho(g_0^u)|\le1$ for every odd $u$.
Writing its eigenvalues as $Q$th roots, the Parseval identity
(K-012d), with $Q$ replacing $M$, again implies
\[
             \sum_{j=0}^{Q/2-1}(a_j-a_{j+Q/2})^2\le1.
\]
A constituent with a nontrivial $\chi$ self-twist has antipodal
pairs because its spectrum is invariant under a nontrivial
two-power rotation.  The same dimension conclusions as in K-013
follow: at most one constituent in a factor can use its unique
pair, that constituent has dimension at most three, and all the
others are lines.  Since $n\ge4$, each factor has at least one
line constituent; this is the only dimension threshold needed here.

Twist each factor by one of its line characters to descend to an
honest $H$-representation.  As in K-013, the product of these three
line characters is a power of $\chi$, so $T$ remains fixed by
$K=\ker\chi$.  In the descended representations the element
$\bar g$ has order $M$.  Ordinary character column orthogonality
on $H$ says that every nonlinear irreducible character vanishes
at $\bar g$.  Hence all nonlinear constituents lie in the unique
paired two-plane, $K$ fixes all the unpaired eigenspaces, and its
remaining action on each paired plane lies in $SL_2$.
An antipodal-free factor is therefore projectively fixed by $K$.
Otherwise a nonzero one-paired-subspace projection gives a nonzero
fixed vector in one of these finite $SL_2$ representations and
forces that action to be trivial.  The connected unit-group
argument of K-012 gives $K=1$ in either case.

For the family of K-015, all hypotheses have been proved there:
criticality and conciseness, $L^0=S$, simple spectra with exactly
one antipodal pair each, a nonzero one-paired-subspace projection,
and diagonal centralizer $S\times\mu_{2^{n-1}}$.
Applying the criterion proves the final assertion.
\end{proof}

Consequently the family issuing from K-004 supplies a genuine
smooth quotient fixed curve, of density one, in every dimension
$n\ge4$.  Its local moving-normal determinant can be used without
an unproved finite-stabilizer assumption.  K-004 supplies the base
dimension four; K-015--K-016 supply all larger dimensions.
This establishes this family of components.  Classification of
the other components and cancellation of their combined leading
coefficients remain the global all-dimensional questions.
